% =====================================================================
%  CONJURA CONJECTURE STATEMENT
%  The polynomial compatibility conjecture.
%  Requires conjura-conjecture.cls in the same directory.
% =====================================================================
\documentclass{conjura-conjecture}

\runninghead{THE POLYNOMIAL COMPATIBILITY CONJECTURE}

\newcommand{\Yc}{\mathcal{Y}}
\newcommand{\Yd}{\hat{\mathcal{Y}}}
\newcommand{\zc}{\hat{0}}

\begin{document}

\cjkicker{CONJURA \textperiodcentered{} OPEN PROBLEM}
\cjtitle{The Polynomial Compatibility Conjecture}
\cjsubtitle{Low-degree, low-influence polynomial distributions over an
  Abelian group}
\cjstatus{Statement: AI-written from Per Austrin, Hao Chung, Kai-Min
  Chung, Shiuan Fu, Yao-Ting Lin, Mohammad Mahmoody, \emph{On the
  Impossibility of Key Agreements from Quantum Random Oracles},
  Cryptology ePrint Archive 2022, not yet checked by a human. The case
  of exponentially small influences, $\delta < |\Yc|^{-d}/d$, is settled
  by the source (Theorem 4.4) and its counterexamples rule out
  $\delta \ge 1/(2d)$; the whole range in between --- in particular the
  inverse-polynomial regime asserted here --- is open.}
\cjcategory{\emph{Idealized Models \& Non-Uniformity}.}

\vspace{0.4em}

\begin{informalconjecture}
Take two probability distributions over real-valued functions on a grid,
each function built from a fixed small number of coordinates at a time
(degree at most $d$), each normalised so that its average squared value
is one, and each coordinate carrying only a tiny share of the total
weight on average (influence at most $\delta$). The conjecture says that
if $\delta$ is small but still only inverse-polynomially small in $d$,
then you can always draw one function from the first distribution and
one from the second and find a single input at which both are nonzero.
This is a purely Fourier-analytic statement, but it was isolated because
it is exactly what is needed to break key agreement in the quantum
random oracle model: the two distributions describe what the oracle
looks like conditioned on the shared key being zero and on it being one,
a common nonzero point is a single oracle consistent with both, and
perfect agreement on the key then makes those two executions contradict
each other. The authors could not prove it. They can prove it only when
the influence bound is exponentially small in the degree, which still
gives an attack, but an exponential-query one. They also give explicit
counterexamples showing the influence bound cannot be relaxed beyond
roughly one over twice the degree, and that dropping either the degree
bound or the influence bound on either of the two distributions makes
the statement false.
\end{informalconjecture}

\section{The setting}\label{sec:setting}

Merkle's puzzles~\cite{Mer74} let two parties who share only a random
oracle agree on a key using $d$ queries while an eavesdropper needs
$\Omega(d^2)$ queries, and a long line of work~\cite{IR89,BMG09,BM17}
showed that this quadratic gap is the best possible in the classical
random oracle model. Quantum access changes the picture: Grover search
breaks Merkle's protocol with $O(d)$ queries, but quantum honest parties
can restore a super-linear gap~\cite{BS08}. Whether parties who compute
quantumly but communicate classically --- the QCCC model --- can achieve
\emph{super-polynomial} security in the quantum random oracle model was
asked by Hosoyamada and Yamakawa~\cite{HY20} and is the motivating
question of the source paper.

The paper reduces that question, for protocols with perfect
completeness, to a statement in Fourier analysis. Following Zhandry's
way of recording quantum queries~\cite{Zha19}, one writes the oracle
register in the Fourier basis; after the two parties make $d$ quantum
queries in total, the joint state of the parties and the oracle is
$d$-sparse there, which is the same thing as saying that the associated
``state polynomial'' on the oracle's truth table has degree at most $d$
and $\ell_2$ norm $1$. A query that the eavesdropper has not yet learned
is one whose variable has small influence, so after an eavesdropper has
learned all $\varepsilon$-heavy queries she is left with two
distributions of exactly the kind in the conjecture: one for the
residual state conditioned on the key being $0$, one conditioned on the
key being $1$. A common point at which a polynomial from each is nonzero
is a single oracle consistent with both executions, and perfect
completeness then says the two executions cannot produce different keys.
So the conjecture is precisely the missing step.

Granting it, the paper obtains a $\poly(d, \log|\Yc|)$-query
\emph{classical} attack on every perfectly complete QCCC key agreement
protocol in the QROM, and, as a corollary, a conditional quantum
black-box separation of QCCC key agreement from one-way functions,
extending it to primitives such as oblivious transfer that imply key
agreement in a black-box way. Unconditionally, the paper proves the
conjecture only in the regime of exponentially small influences,
$\delta < |\Yc|^{-d}/d$, which yields an attack of $2^{O(dm)} \cdot d^2$
queries for oracles with $m$-bit outputs. Its counterexamples show the
hypotheses cannot be dropped: both distributions need the degree bound,
both need the influence bound, and $\delta$ must be below $1/(2d)$.

The statement is a cousin of the Aaronson--Ambainis
conjecture~\cite{AA14}, whose contrapositive also asserts that a
low-degree polynomial with polynomially small influences must be large
somewhere. Both are known when the influences are exponentially
small~\cite{DFKO06}, and the degenerate case of no communication and
perfect completeness is settled by~\cite{OSSS05}; but the paper reports
that the two conjectures do not appear to be directly comparable, and
that their cryptographic consequences are incomparable too.

Progress to date. Theorem 4.4 proves the conjecture for
$\delta < |\Yc|^{-d}/d$ by fixing a maximum-degree term of $f$ and
showing that a $g$ with small influences has, in expectation over a
random assignment to the remaining variables, a constant term dominating
all $|\Yc|^{d}$ non-constant terms of the restriction; the $|\Yc|^{-d}$
loss comes from that count (in the Boolean case $\Yc = \mathbb{Z}_2$
this is the $2^{-d}$ of the paper's overview). The proof in fact uses no
influence condition on $F$ and no degree bound on $G$. Theorem 5.6
(Appendix A) shows it suffices to treat real-valued rather than
complex-valued functions. Lemma 4.8 (Section 8) shows that the
conjecture for one constant-size Abelian group already gives
$\poly(d, \log|\Yc'|)$-query attacks for oracles with range any finite
Abelian group $\Yc'$, so the choice of group in the conjecture is not
the obstacle.

\section{Notation and parameters}\label{sec:notation}

Throughout, $\Yc$ is a finite Abelian group, written additively, and
$\Yd$ is its dual group, i.e.\ the set of characters
$\chi \colon \Yc \to \{\omega \in \mathbb{C} : |\omega| = 1\}$ satisfying
$\chi(y + y') = \chi(y)\chi(y')$. The trivial character, $\chi(y) = 1$
for all $y \in \Yc$, is denoted $\zc$. One has $|\Yd| = |\Yc|$.

For $N \in \mathbb{N}$ write $[N] = \{1, \dots, N\}$. A tuple
$\chi = (\chi_1, \dots, \chi_N) \in \Yd^N$ is viewed as the function
$\chi(x) = \prod_{i=1}^{N} \chi_i(x_i)$ on
$x = (x_1, \dots, x_N) \in \Yc^N$. These $|\Yc|^N$ functions form an
orthonormal basis for the complex-valued functions on $\Yc^N$ under the
inner product
$\langle f, g \rangle = \mathbb{E}_{x}[f(x)\overline{g(x)}]$, where $x$
is uniform on $\Yc^N$. Consequently every $f \colon \Yc^N \to \mathbb{R}$
has a unique Fourier expansion
\[
  f(x) \;=\; \sum_{\chi \in \Yd^N} \hat{f}(\chi)\, \chi(x),
  \qquad \hat{f}(\chi) \in \mathbb{C}.
\]
The $\ell_2$ norm is
\[
  \|f\|_2 \;=\; \bigl(\mathbb{E}_{x}[|f(x)|^2]\bigr)^{1/2}
  \;=\; \Bigl(\sum_{\chi} |\hat{f}(\chi)|^2\Bigr)^{1/2},
\]
again with $x$ uniform on $\Yc^N$.

For a probability distribution $F$ over functions
$\Yc^N \to \mathbb{R}$, $\operatorname{supp}(F)$ denotes its support,
and $f \sample F$ denotes a sample. The parameters are the degree bound
$d \in \mathbb{N}$, the number of variables $N \in \mathbb{N}$, and the
influence bound $\delta$.

\begin{itemize}
\item $\Yc$: a finite Abelian group; the range of the random oracle in
  the cryptographic application. The conjecture asserts that some such
  group works; the Boolean case $\Yc = \mathbb{Z}_2$ is the special case
  stated as Conjecture 1.2 in the paper.
\item $N$: number of variables, i.e.\ the size of the oracle's domain.
  Universally quantified; the influence bound must not depend on it.
\item $d$: degree bound on the polynomials; in the application, the
  total number of quantum oracle queries made by the two honest parties.
\item $\delta(d)$: bound on the average influence of each variable. The
  conjecture requires $\delta(d)$ to be at least inverse polynomial in
  $d$; the paper's counterexamples force $\delta(d) < 1/(2d)$.
\item $F, G$: the two distributions over polynomials. In the application
  they describe the oracle's state conditioned on the agreed key bit
  being $0$ and being $1$ respectively.
\end{itemize}

\section{The conjecture}\label{sec:conjectures}

\begin{definition}[degree and influence]\label{def:deginf}
For $\chi = (\chi_1, \dots, \chi_N) \in \Yd^N$ let
$\deg(\chi) = \bigl|\{ i \in [N] : \chi_i \ne \zc \}\bigr|$. The
\emph{degree} of $f \colon \Yc^N \to \mathbb{R}$ is
\[
  \deg(f) \;=\; \max\{ \deg(\chi) \;:\; \hat{f}(\chi) \ne 0 \},
\]
with $\deg(f) = 0$ for $f$ constant. For $i \in [N]$, the
\emph{influence} of the $i$-th variable on $f$ is
\[
  \operatorname{Inf}_i(f) \;=\;
  \sum_{\chi \,:\, \chi_i \ne \zc} \bigl|\hat{f}(\chi)\bigr|^2 .
\]
For $\Yc = \mathbb{Z}_2$, identifying $\Yc^N$ with $\{\pm 1\}^N$, these
are the usual notions for multilinear polynomials: writing
$f = \sum_{S \subseteq [N]} \alpha_S \prod_{i \in S} x_i$ with
$\alpha_S \in \mathbb{R}$, one has
\[
  \deg(f) = \max_{\alpha_S \ne 0} |S|, \qquad
  \operatorname{Inf}_i(f) = \sum_{S \ni i} \alpha_S^2, \qquad
  \|f\|_2^2 = \sum_S \alpha_S^2 .
\]
\end{definition}

\begin{conjecture}[polynomial compatibility]\label{conj:main}
There exist a finite Abelian group $\Yc$, a constant $c > 0$, and a
function $\delta \colon \mathbb{N} \to (0,1]$ with
$\delta(d) \ge d^{-c}$ for all $d \in \mathbb{N}$, such that the
following holds for all $d, N \in \mathbb{N}$. Let $F$ and $G$ be
probability distributions over functions $\Yc^N \to \mathbb{R}$ such
that:
\begin{itemize}
  \item \emph{unit $\ell_2$ norm:} $\|f\|_2 = 1$ for every
    $f \in \operatorname{supp}(F)$ and $\|g\|_2 = 1$ for every
    $g \in \operatorname{supp}(G)$;
  \item \emph{degree at most $d$:} $\deg(f) \le d$ for every
    $f \in \operatorname{supp}(F)$ and $\deg(g) \le d$ for every
    $g \in \operatorname{supp}(G)$;
  \item \emph{$\delta$-influences on average:} for every $i \in [N]$,
    \begin{align*}
      \mathbb{E}_{f \sample F}\bigl[\operatorname{Inf}_i(f)\bigr]
        &\le \delta(d) \quad\text{and}\\
      \mathbb{E}_{g \sample G}\bigl[\operatorname{Inf}_i(g)\bigr]
        &\le \delta(d).
    \end{align*}
\end{itemize}
Then there exist $f \in \operatorname{supp}(F)$,
$g \in \operatorname{supp}(G)$ and $x \in \Yc^N$ such that
$f(x) \cdot g(x) \ne 0$.
\end{conjecture}

\begin{remark}[which group, and how it is quantified]\label{rem:group}
Conjecture~\ref{conj:main} is the \emph{existential} form: it asks only
that some finite Abelian group $\Yc$ admit an inverse-polynomial
$\delta$. This is the form taken by the source's Conjectures 4.3 and
5.5, and it is the form its Theorems 4.5 and 6.3 use; by Lemma 4.8 one
constant-size group already suffices for attacks against oracles whose
range is any finite Abelian group, so the choice of group is not where
the difficulty lies. The paper's introductory Conjecture 1.2 fixes the
Boolean case $\Yc = \mathbb{Z}_2$; that version is strictly stronger and
also open, and it must not be presented as the paper's main conjecture.
\end{remark}

\begin{remark}[how to read the hypotheses]\label{rem:reading}
Three points are easy to get wrong. First, the influence bound is on the
\emph{average} over the distribution,
$\mathbb{E}_{f \sample F}[\operatorname{Inf}_i(f)]$, and not on each
polynomial in the support; imposing it pointwise strengthens the
hypothesis and so weakens the conjecture. Second, $N$ is universally
quantified and $\delta$ depends on $d$ alone: the source's Conjecture
5.5 leaves $N$ free in the text, and the universal reading is the one
made explicit in the equivalent Conjecture 4.3 (``for any
$d, N \in \mathbb{N}$''). Third, the paper's introduction writes the
$\ell_2$ norm without the square root; since the constraint is that the
norm equals $1$ this makes no difference, and the standard normalisation
of Section~\ref{sec:notation} is used here.
\end{remark}

\begin{remark}[the window that is open]\label{rem:window}
The source proves Conjecture~\ref{conj:main} (Theorem 4.4) whenever
$\delta < |\Yc|^{-d}/d$, i.e.\ for exponentially small influences, and
shows (Theorem 5.6) that this polynomial formulation is equivalent to
the quantum-state formulation, its Conjecture 4.3. Appendix B rules out
$\delta \ge 1/(2d)$ and shows that the degree and influence hypotheses
are needed on both distributions. Everything between $|\Yc|^{-d}/d$ and
$1/(2d)$ is open, and any $\delta$ that is $1/\poly(d,\log|\Yc|)$
already suffices for a polynomial-query attack.
\end{remark}

\begin{remark}[what settling it would give]\label{rem:consequence}
The conjecture is the whole gap between what is known and a full answer
to the question of whether quantum parties talking over a classical
channel can get super-polynomial security from a random oracle, at least
in the perfectly complete case. Proving it gives an unconditional
polynomial-query classical attack on every such protocol and a quantum
black-box separation of QCCC key agreement --- hence oblivious transfer
and public-key encryption with classical inputs and communication ---
from one-way functions. Refuting it would not immediately give a secure
protocol, but it would kill the only known route to one of the two main
barriers in this area, and it would be a genuinely new phenomenon in the
Fourier analysis of low-influence low-degree functions, a regime where
the analogous Aaronson--Ambainis conjecture~\cite{AA14} is widely
believed.
\end{remark}

\section{Bibliography}

\begin{conjurabibliography}{99}

\bibitem[AA14]{AA14}
S. Aaronson, A. Ambainis.
\emph{The Need for Structure in Quantum Speedups}.
2014. No venue printed in the source bibliography.

\bibitem[BM17]{BM17}
B. Barak, M. Mahmoody.
\emph{Merkle's Key Agreement Protocol is Optimal: An $O(n^2)$ Attack on
Any Key Agreement from Random Oracles}.
Journal of Cryptology, 30(3), 2017.

\bibitem[BMG09]{BMG09}
B. Barak, M. Mahmoody-Ghidary.
\emph{Merkle Puzzles are Optimal --- An $O(n^2)$-Query Attack on Any Key
Exchange from a Random Oracle}.
Advances in Cryptology --- CRYPTO 2009, LNCS 5677, pages 374--390,
Springer, 2009.

\bibitem[BS08]{BS08}
G. Brassard, L. Salvail.
\emph{Quantum Merkle Puzzles}.
International Conference on Quantum, Nano and Micro Technologies
(ICQNM), pages 76--79, IEEE Computer Society, 2008.

\bibitem[DFKO06]{DFKO06}
I. Dinur, E. Friedgut, G. Kindler, R. O'Donnell.
\emph{On the Fourier Tails of Bounded Functions over the Discrete Cube}.
Proceedings of the Thirty-Eighth Annual ACM Symposium on Theory of
Computing, pages 437--446, 2006.

\bibitem[HY20]{HY20}
A. Hosoyamada, T. Yamakawa.
\emph{Finding Collisions in a Quantum World: Quantum Black-Box
Separation of Collision-Resistance and One-Wayness}.
International Conference on the Theory and Application of Cryptology and
Information Security, pages 3--32, Springer, 2020.

\bibitem[IR89]{IR89}
R. Impagliazzo, S. Rudich.
\emph{Limits on the Provable Consequences of One-Way Permutations}.
Proceedings of the 21st Annual ACM Symposium on Theory of Computing
(STOC), pages 44--61, ACM Press, 1989.

\bibitem[Mer74]{Mer74}
R. Merkle.
\emph{C.s.\ 244 Project Proposal}.
1974. Facsimile available at \url{http://www.merkle.com/1974}.

\bibitem[OSSS05]{OSSS05}
R. O'Donnell, M. Saks, O. Schramm, R. A. Servedio.
\emph{Every Decision Tree Has an Influential Variable}.
46th Annual IEEE Symposium on Foundations of Computer Science
(FOCS'05), pages 31--39, IEEE, 2005.

\bibitem[Zha19]{Zha19}
M. Zhandry.
\emph{How to Record Quantum Queries, and Applications to Quantum
Indifferentiability}.
Annual International Cryptology Conference, pages 239--268, Springer,
2019.

\end{conjurabibliography}

\end{document}
